\documentclass[german]{cv}

\hypersetup{
    pdftitle={Bewerbungsschreiben von Jeongjun Ahn},
    pdfauthor={Jeongjun Ahn}
}
\cvname{Jeongjun Ahn}

% Letters need visible paragraph separation (the class sets \parskip to 0):
\setlength{\parskip}{0.5em}

\begin{document}

    \pagestyle{empty} % Brief ohne CV-Fußzeile

    % Absenderblock:
    Jeongjun Ahn\\
    \platzhalter{Straße und Hausnummer}\\
    \platzhalter{Postleitzahl und Ort}\\
    \platzhalter{Telefonnummer}\\
    \hrefWithoutArrow{mailto:tyyyu1268@gmail.com}{\color{black}tyyyu1268@gmail.com}

    \vspace{0.8 cm}

    % Empfänger:
    \platzhalter{Offizielle Bezeichnung des Unternehmens}\\
    \platzhalter{Abteilung, falls bekannt}\\
    \platzhalter{Herrn/Frau Name der Ansprechperson, falls vorhanden}\\
    \platzhalter{Straße und Hausnummer}\\
    \platzhalter{Postleitzahl und Ort}

    \vspace{0.2 cm}

    \begin{flushright}
        Leipzig, \platzhalter{TT. Monat JJJJ}
    \end{flushright}

    \vspace{0.2 cm}

    \textbf{Bewerbung als \platzhalter{genaue Stellenbezeichnung}}

    \vspace{0.4 cm}

    % Einstieg: Motivation für die Bewerbung und konkretes Interesse am Unternehmen.
    % Persönliche Vorstellung: warum ich zur Stelle passe. Stärken und Fähigkeiten mit Beispielen belegen.
    % Verfügbarkeit
    % Schluss: deutlich machen, dass ich mich über eine Einladung zum Vorstellungsgespräch freue.

    \platzhalter{Sehr geehrte Frau .../Sehr geehrter Herr ...,}\\
    \platzhalter{ohne Ansprechperson: Sehr geehrte Damen und Herren,}

    \platzhalter{Einstieg: Motivation für die Bewerbung und konkretes Interesse am Unternehmen. Kein Standardsatz - Bezug zur Stelle oder zum Produkt herstellen. Nach der Anrede klein weiterschreiben.}

    \platzhalter{Praktische Erfahrung: Studium, Bachelorarbeit und eigene Projekte, mit konkreten Beispielen belegt.}

    \platzhalter{Arbeitsweise und Stärken: wie Sie Probleme angehen und worauf Sie bei Software besonders Wert legen.}

    \platzhalter{Bezug zum Tech-Stack des Unternehmens und was Sie dort lernen möchten.}

    \platzhalter{Verfügbarkeit und Schluss: möglicher Eintrittstermin; deutlich machen, dass Sie sich über eine Einladung zum Vorstellungsgespräch freuen.}

    \vspace{0.4 cm}

    Mit freundlichen Grüßen

    % \includegraphics[height=1.5 cm]{../images/signature_jj.png}
    \platzhalter{Unterschrift}

    Jeongjun Ahn

    \vspace{0.3 cm}

    \textbf{Anlagen}\\
    Lebenslauf\\
    \platzhalter{Weitere Anlagen, z.\,B. Empfehlungsschreiben, Zeugnisse}

\end{document}
